\documentclass[instructions.tex]{subfiles}
\begin{document}
 
\chapter{Travailler au salut avec de grands soins}
 
\primo Puisqu'il s'agit de tout pour l'éternité dans l'affaire du salut, pourquoi y travaille-t-on avec si peu de soin ? Le ciel à gagner ou à perdre mérite-t-il moins d'attention que les choes de la terre ? Ne seriez-vous pas un insensé si, ayant un procès où il s'agit de votre fortune et de votre vie, vous ne pensiez qu'à vous divertir à la veille de le perdre ? Vous êtes bien plus insensé si vous négligez votre salut, puisqu'il s'y agit de tout gagner ou de tout perdre pour toujours.

Que l'aveuglement des hommes est déplorable! Le plupart trouvent du temps pour des occupations et des recherches infructueuses, et ils n'en trouvent point pour vivre en chrétiens et pour se sauver, dit saint Paulin\footnote{Vacat tibi ut philosophus sis, non vacat ut christianus sis. S.Paulin}. Ils ont soin que leurs maisons soient rétablies, leurs terres cultivées, leurs revenus assurés. Que d'applications pour conserver leurs fonds, leur santé, leur beauté, leurs vêtements! mais ont-ils la même ardeur pour ce qui regarde le salut de leur âme ? Quand ce serait l'âme de leur ennemi ou l'âme d'une bête, la traiteraient-ils plus mal ? On dirait qu'ils n'ont une âme que pour la perdre. Hélas! à quoi aboutissent nos travaux et toute ce que nous faisons au monde, si nous n'y faisons ce que nous devons uniquement y faire ? 

\secundo Le salut n'est pas une affaire qui réussisse par hasard; il demande de grands soins, il demande même tous nos soins. C'est pour cela que Jésus-Christ a si souvent répété ces paroles: \og Faites vos efforts pour entrer par la porte étroite qui conduit à la vie\fg. Peut-on, en effet, prendre top de précautions et de soins pour ne pas risquer la perte d'une âme immortelle ? Or, quelles précautions devriez-vous prendre ? 

Apprenez-le du fils de Dieu. Voyez les supplices et la mort qu'il a soufferts pour sauver votre âme. Apprenez-le des martyrs, qui ont mieux aimé souffrir les tourmens et les tortures, que de perdre la leur en perdant la foi; apprenez-le des saints pénitents, qui ont crucifié leur chair, qui ont vécu dans le détachement, pour ne s'occuper que du soin de leur âme; apprenez-le de tant de personnes qui ont fait à Dieu les plus grands sacrifices, qui ont résisté aux puissances mêmes. Tel fut un grand pape qui, refusant à un roi une chose qu'il ne pouvait accorder sans préjudicier à son âme, lui dit ces mémorables paroles:\og si j'avais deux âmes, j'en pourrais risquer une pour vous satisfaire; mais je n'en ai qu'une, et rien au monde ne peut m'engager à risquer de la perdre\fg. Apprenez du démon même le soin que vous devez avoir de votre âme. N'est-il pas étonnant qu'il prenne plus de précautions pour la perdre que nous n'en prenons pour la sauver ? Quelle fureur, dit Salvien, de n'estimer pas votre âme digne de vos soins, pendant que le démon la met à un si haut prix qu'il fait tous ses efforts pour l'avoir! Quelle honte pour vous, puisqu'en négligeant le soin de votre âme, vous l'estimez beaucoup moins que le démon ne l'estime!

\section{Résolutions}

Trois ennemis peuvent m'empêcher de travailler à mon salut avec soin: le monde, le démon, et mes inclinations mauvaises. 

\primo Je vaincrai le monde, en ne me produisant plus désormais dans ses divertissements dangereux; n'ayant même de rapports avec lui, que ceux que la nécessité, l'utilité ou la bien-séance exigeront, et toujours avec de sérieuses précautions.

\secundo Je combattrai le démon par la confiance en Dieu, l'appelant à mon secours, aussitôt que je m'apperceverai que je suis tenté.

\tertio Je ferais sans cesse la guerre aux mauvais penchans qui dominent en moi: si c'est l'orgueil, je l'abbatrai par l'humilité; si c'est l'attachement aux biens terrestres, je ferai l'aumône autant que mes facultés me le permettront, rejetant les vains désirs de fortune et d'aggrandissement; si c'est l'amour du plaisir, j'y opposerai une mortification continuelle, qui s'étendra sur ma nourriture, mes sens, ma langue, mes pensées, retranchant avec un soin extrême tout ce qui pourrait enflammer une inclination si dangereuse. 

\prayer{O mon Dieu! je vous confie ces résolutions, rappelez-les souvent à ma mémoire, et faites que j'y sois fidèle le reste de mes jours.}

\end{document}
